% GENERATED FILE. Edit canonical lessons, then run npm run generate:books.
% canonical-source-hash: fnv1a64:4fa74cd6f06f7eee
% canonical-lessons: UR-C402-masafat, UR-C402-rahnuma, UR-C402-thakana, UR-C402-sulana, UR-C402-jagana

\chapter{Long Journeys}
\label{ch:ur-long-journeys}
\hlchaptermodality{402}

\begin{chapteropening}
\textbf{By the end of this chapter:} \emph{I can talk about a long journey: the distance, the guide and a tired child.}

Complete the last lesson of chapter 402: I can talk about a long journey: the distance, the guide and a tired child.
\end{chapteropening}

\section[\texorpdfstring{\ur{مسافت} (masāfat)}{مسافت (masāfat)}]{\ur{مسافت} (masāfat) — a distance travelled; a long way}

\label{lesson:UR-C402-masafat}



\begin{quote}
Before the new one: say the Urdu for a bouquet, then the Urdu for to wave.
\end{quote}

\subsection*{You'll want to know: \ur{مسافت}}
\textbf{\ur{مسافت}} — \emph{masāfat} — \textquotedblleft{}a distance travelled; a long way\textquotedblright{}.

\begin{quote}
\ur{پہاڑوں} \ur{تک} \ur{کی} \ur{مسافت} \ur{لمبی} \ur{ہے۔}
\end{quote}

\begin{quote}
\emph{pahāṛoṅ tak kī masāfat lambī hai.}
\end{quote}

\begin{quote}
It is a long way to the mountains.
\end{quote}

Feminine.

\begin{cousinweb}
\textbf{\ur{مسافت}} (\emph{masāfat}) is usually traced, through Persian, to Arabic \emph{masāfa}, \textquotedblleft{}a distance\textquotedblright{}; it is a cousin of \textbf{\ur{مسافر}} and \textbf{\ur{سفر}}.
\end{cousinweb}

\begin{grammarlens}[title={What you've built}]
One more word for getting around.
\end{grammarlens}

\subsection*{Your turn}
\begin{itemize}
\raggedright
  \item \emph{Say it:} \emph{masāfat}
  \item \emph{Say it:} \emph{masāfat}, once more
  \item \emph{Say it:} \emph{pahāṛoṅ tak kī masāfat lambī hai.} — It is a long way to the mountains.
  \item \emph{Recall:} say the Urdu for a bouquet, then the Urdu for to wave, then say \emph{masāfat} again
\end{itemize}

\begin{tcolorbox}[breakable,colback=teal!4,colframe=teal!35!black,title={Before you move on}]
What does \ur{مسافت} mean? (\textquotedblleft{}A distance travelled; a long way\textquotedblright{}.) Say it once more.
\end{tcolorbox}
\section[\texorpdfstring{\ur{رہنما} (rahnumā)}{رہنما (rahnumā)}]{\ur{رہنما} (rahnumā) — a guide; a leader}

\label{lesson:UR-C402-rahnuma}



\begin{quote}
Before the new one: say the Urdu for to wave, then the Urdu for a distance travelled.
\end{quote}

\subsection*{You'll want to know: \ur{رہنما}}
\textbf{\ur{رہنما}} — \emph{rahnumā} — \textquotedblleft{}a guide; a leader\textquotedblright{}.

\begin{quote}
\ur{پہاڑوں} \ur{میں} \ur{ایک} \ur{رہنما} \ur{ہمارے} \ur{ساتھ} \ur{تھا۔}
\end{quote}

\begin{quote}
\emph{pahāṛoṅ meṅ ek rahnumā hamāre sāth thā.}
\end{quote}

\begin{quote}
In the mountains a guide went with us.
\end{quote}

Masculine or feminine, by the person meant.

\begin{cousinweb}
\textbf{\ur{رہنما}} (\emph{rahnumā}) is usually traced to Persian \emph{rahnumā}, \textquotedblleft{}one who shows the way\textquotedblright{}, built from \emph{rah}, \textquotedblleft{}a way\textquotedblright{} (a short form of \textbf{\ur{راہ}}), and \emph{numā}, \textquotedblleft{}showing\textquotedblright{}.
\end{cousinweb}

\begin{grammarlens}[title={What you've built}]
One more word for getting around.
\end{grammarlens}

\subsection*{Your turn}
\begin{itemize}
\raggedright
  \item \emph{Say it:} \emph{rahnumā}
  \item \emph{Say it:} \emph{rahnumā}, once more
  \item \emph{Say it:} \emph{pahāṛoṅ meṅ ek rahnumā hamāre sāth thā.} — In the mountains a guide went with us.
  \item \emph{Recall:} say the Urdu for to wave, then the Urdu for a distance travelled, then say \emph{rahnumā} again
\end{itemize}

\begin{tcolorbox}[breakable,colback=teal!4,colframe=teal!35!black,title={Before you move on}]
What does \ur{رہنما} mean? (\textquotedblleft{}A guide; a leader\textquotedblright{}.) Say it once more.
\end{tcolorbox}
\section[\texorpdfstring{\ur{تھکانا} (thakānā)}{تھکانا (thakānā)}]{\ur{تھکانا} (thakānā) — to tire (someone) out}

\label{lesson:UR-C402-thakana}



\begin{quote}
Before the new one: say the Urdu for a distance travelled, then the Urdu for a guide.
\end{quote}

\subsection*{You'll want to know: \ur{تھکانا}}
\textbf{\ur{تھکانا}} — \emph{thakānā} — \textquotedblleft{}to tire (someone) out\textquotedblright{}.

\begin{quote}
\ur{لمبے} \ur{سفر} \ur{نے} \ur{سارا} \ur{کو} \ur{تھکا} \ur{دیا۔}
\end{quote}

\begin{quote}
\emph{lambe safar ne sārā ko thakā diyā.}
\end{quote}

\begin{quote}
The long journey wore Sara out.
\end{quote}

Past, with \ur{نے} (he or she …): \ur{اس} \ur{نے} \ur{گھوڑے} \ur{کو} \ur{تھکایا}, \emph{us ne ghoṛe ko thakāyā}.

\begin{cousinweb}
\textbf{\ur{تھکانا}} (\emph{thakānā}) is built on \textbf{\ur{تھکنا}}, \textquotedblleft{}to get tired\textquotedblright{}, which comes from \emph{sthakk}, \textquotedblleft{}stop, halt\textquotedblright{}, a form CDIAL reconstructs (CDIAL 13737): this is the form for tiring someone else.
\end{cousinweb}

\begin{grammarlens}[title={What you've built}]
One more word for getting around.
\end{grammarlens}

\subsection*{Your turn}
\begin{itemize}
\raggedright
  \item \emph{Say it:} \emph{thakānā}
  \item \emph{Say it:} \emph{thakānā}, once more
  \item \emph{Say it:} \emph{lambe safar ne sārā ko thakā diyā.} — The long journey wore Sara out.
  \item \emph{Recall:} say the Urdu for a distance travelled, then the Urdu for a guide, then say \emph{thakānā} again
\end{itemize}

\begin{tcolorbox}[breakable,colback=teal!4,colframe=teal!35!black,title={Before you move on}]
What does \ur{تھکانا} mean? (\textquotedblleft{}To tire (someone) out\textquotedblright{}.) Say it once more.
\end{tcolorbox}
\section[\texorpdfstring{\ur{سلانا} (sulānā)}{سلانا (sulānā)}]{\ur{سلانا} (sulānā) — to put (someone) to sleep}

\label{lesson:UR-C402-sulana}



\begin{quote}
Before the new one: say the Urdu for a guide, then the Urdu for to tire (someone) out.
\end{quote}

\subsection*{You'll want to know: \ur{سلانا}}
\textbf{\ur{سلانا}} — \emph{sulānā} — \textquotedblleft{}to put (someone) to sleep\textquotedblright{}.

\begin{quote}
\ur{ماں} \ur{نے} \ur{گاڑی} \ur{میں} \ur{بچے} \ur{کو} \ur{سلایا۔}
\end{quote}

\begin{quote}
\emph{māṅ ne gāṛī meṅ bachche ko sulāyā.}
\end{quote}

\begin{quote}
The mother got the child off to sleep in the car.
\end{quote}

Past, with \ur{نے} (he or she …): \ur{اس} \ur{نے} \ur{بچے} \ur{کو} \ur{سلایا}, \emph{us ne bachche ko sulāyā}.

\begin{cousinweb}
\textbf{\ur{سلانا}} (\emph{sulānā}) is inherited: CDIAL lists the Hindi cousin \emph{sulānā}, \textquotedblleft{}lulls to sleep\textquotedblright{}, under Sanskrit \emph{svāpáyati} (CDIAL 13928); it is the making-sleep partner of \textbf{\ur{سونا}}.
\end{cousinweb}

\begin{grammarlens}[title={What you've built}]
One more word for getting around.
\end{grammarlens}

\subsection*{Your turn}
\begin{itemize}
\raggedright
  \item \emph{Say it:} \emph{sulānā}
  \item \emph{Say it:} \emph{sulānā}, once more
  \item \emph{Say it:} \emph{māṅ ne gāṛī meṅ bachche ko sulāyā.} — The mother got the child off to sleep in the car.
  \item \emph{Recall:} say the Urdu for a guide, then the Urdu for to tire (someone) out, then say \emph{sulānā} again
\end{itemize}

\begin{tcolorbox}[breakable,colback=teal!4,colframe=teal!35!black,title={Before you move on}]
What does \ur{سلانا} mean? (\textquotedblleft{}To put (someone) to sleep\textquotedblright{}.) Say it once more.
\end{tcolorbox}
\section[\texorpdfstring{\ur{جگانا} (jagānā)}{جگانا (jagānā)}]{\ur{جگانا} (jagānā) — to wake (someone) up}

\label{lesson:UR-C402-jagana}



\begin{quote}
Before the new one: say the Urdu for to tire (someone) out, then the Urdu for to put (someone) to sleep.
\end{quote}

\subsection*{You'll want to know: \ur{جگانا}}
\textbf{\ur{جگانا}} — \emph{jagānā} — \textquotedblleft{}to wake (someone) up\textquotedblright{}.

\begin{quote}
\ur{کنڈکٹر} \ur{نے} \ur{صبح} \ur{سب} \ur{مسافروں} \ur{کو} \ur{جگایا۔}
\end{quote}

\begin{quote}
\emph{kanḍakṭar ne subah sab musāfiroṅ ko jagāyā.}
\end{quote}

\begin{quote}
In the morning the conductor woke all the passengers.
\end{quote}

Past, with \ur{نے} (he or she …): \ur{اس} \ur{نے} \ur{سارا} \ur{کو} \ur{جگایا}, \emph{us ne sārā ko jagāyā}.

\begin{cousinweb}
\textbf{\ur{جگانا}} (\emph{jagānā}) is inherited: CDIAL lists Hindi \emph{jagānā} beside \emph{jāgnā}, \textquotedblleft{}to wake\textquotedblright{}, under Sanskrit \emph{jāgrati}, \textquotedblleft{}are awake\textquotedblright{} (CDIAL 5175).
\end{cousinweb}

\begin{grammarlens}[title={What you've built}]
One more word for getting around.
\end{grammarlens}

\subsection*{Your turn}
\begin{itemize}
\raggedright
  \item \emph{Say it:} \emph{jagānā}
  \item \emph{Say it:} \emph{jagānā}, once more
  \item \emph{Say it:} \emph{kanḍakṭar ne subah sab musāfiroṅ ko jagāyā.} — In the morning the conductor woke all the passengers.
  \item \emph{Recall:} say the Urdu for to tire (someone) out, then the Urdu for to put (someone) to sleep, then say \emph{jagānā} again
\end{itemize}

\begin{tcolorbox}[breakable,colback=teal!4,colframe=teal!35!black,title={Before you move on}]
What does \ur{جگانا} mean? (\textquotedblleft{}To wake (someone) up\textquotedblright{}.) Say it once more.
\end{tcolorbox}
